% Originally: content/week-12.tex, \section{Pullback of forms} (Corsin Week 12)


\section{Pullback of forms}
\label{sec:pullback_of_forms}

% Source: Corsin Nick/Class Notes/Week 12.pdf, p. 5
\begin{definition}[Pullback of a form]
\label{def:pullback_of_forms}
Let $U \subseteq \mathbb{R}^m$, $V \subseteq \mathbb{R}^n$ be open, $F : U \to V$ a smooth
map, and $\omega \in \Omega^p(V)$ a differential $p$-form on $V$. The \newterm{pullback}
\germanfor{Zur\"uckziehung} of $\omega$ under $F$ is the differential $p$-form on $U$
defined by
\[(F^*\omega)_x(v_1,\ldots,v_p) := \omega_{F(x)}\big(DF_x(v_1),\ldots,DF_x(v_p)\big).\]
\end{definition}

% Moved out of the typeset text (the book names tutors only in source citations):
% Corsin also draws a commutative diagram TU --dF--> TV --omega--> R over U --F--> V, with
% tangent-bundle projections pi_U, pi_V, and labels it "not exam relevant". Omitted here for that
% reason; the formula above is the operative content.

The pullback satisfies two properties that make it the right tool for transporting forms
along smooth maps:
\begin{enumerate}[label=\textbf{(\alph*)}]
  \item $F^*(\omega\wedge\theta) = F^*\omega \wedge F^*\theta$;
  \item $F^*(d\omega) = d(F^*\omega)$ (the pullback commutes with the exterior derivative).
\end{enumerate}

\begin{ainote}
This is precisely the operation used, informally, in the reparametrization-invariance
motivation for \cref{def:antisymmetric_k_linear_form}: writing $\phi^*\omega$ and
$(\phi\circ\Psi)^*\omega$ for the two sides of that calculation makes it read as $\int_D
\phi^*\omega = \int_V \Psi^*(\phi^*\omega) = \int_V (\phi\circ\Psi)^*\omega$. The first step is
the change of variables formula for a top-degree form, $\int_D\eta = \int_V\Psi^*\eta$ when
$\det D\Psi > 0$ (\cref{ex:change_of_variables_differential_forms}), applied to
$\eta := \phi^*\omega$; the second is the fact that pullback along a composition is the
composition of pullbacks, $(\phi\circ\Psi)^* = \Psi^*\circ\phi^*$, which is the chain rule.
% Correction: Claude Opus 5.5 (Extra) --- attributed the calculation to "property (a)", which is
% the wedge-product rule and plays no part in it.
\end{ainote}

% Source: Jérôme Paschoud/Notizen/Woche 12.1 Differentialformen.pdf, p. 7
\begin{example}[Computing a pullback]
\label{ex:computing_pullback}
Let $\omega = x_3 \,dx_3 + x_1 x_2 \,dx_2$ be a $1$-form on $\mathbb{R}^3$.
Consider the map $F : \mathbb{R}^2 \to \mathbb{R}^3$ given by
\[F(u,v) = \begin{pmatrix} u+v \\ e^u \\ u^2 - v \end{pmatrix}.\]
We want to compute the pullback $F^*\omega$, which will be a $1$-form on $\mathbb{R}^2$.
Using the properties of the pullback, we have:
\begin{align*}
F^*\omega &= F^*(x_3 \,dx_3 + x_1 x_2 \,dx_2) \\
&= F^*(x_3) F^*(dx_3) + F^*(x_1 x_2) F^*(dx_2) \\
&= (x_3 \circ F) d(F^*(x_3)) + (x_1 x_2 \circ F) d(F^*(x_2)) \\
&= (u^2 - v) d(u^2 - v) + (u+v)e^u d(e^u) \\
&= (u^2 - v)(2u \,du - dv) + (u+v)e^u (e^u \,du) \\
&= \big(2u(u^2 - v) + (u+v)e^{2u}\big) \,du - (u^2 - v) \,dv.
\end{align*}
\end{example}

\begin{aiexercise}[The pullback of the area form is the Jacobian]
% Generator: Claude Opus 5 (High)
\label{ex:ai_pullback_area_form}
Let $\Phi(r,\theta) := (r\cos\theta,\ r\sin\theta)$ be the polar coordinate map, and write
$(x,y)$ for the coordinates on the target.
\begin{enumerate}[label=\textbf{(\alph*)}]
  \item Compute $\Phi^*(dx)$ and $\Phi^*(dy)$, and wedge them together to obtain
    $\Phi^*(dx\wedge dy)$. Compare the answer with the polar entry of
    \cref{rem:coordinate_jacobians_table}.
  \item \Cref{ex:change_of_variables_differential_forms} proves the general identity behind part \textbf{(a)}, namely that a
    pullback of a top-degree form produces the factor $\det\Jac\Phi$, and it assumes
    $\det\Jac\Phi > 0$ throughout. \Cref{thm:change_of_variables}, by contrast, carries
    $\lvert\det\Jac\Phi\rvert$ and needs no such hypothesis. Explain what the two are measuring
    differently, and which of them notices that $\Phi(u_1,u_2) := (u_2,u_1)$ reverses
    orientation.
  \item Using part \textbf{(b)}, say what would go wrong in \cref{ex:change_of_variables_differential_forms} if its positivity
    hypothesis were dropped, and why the change of variables formula is untroubled by the same
    map.
\end{enumerate}
\exsol{sol:ai_pullback_area_form}
\end{aiexercise}
